\chapter{Strabismus}

\section*{Definition}
Strabismus (squint) is a misalignment of the visual axes such that both eyes do not simultaneously fixate on the same point of regard; deviation may be constant or intermittent, comitant (same angle in all gazes) or incomitant (varying with direction of gaze).

\section*{Pathophysiology}
Imbalance in the extraocular muscle tone, abnormal neuromuscular control, or disruption of the vergence and fusion systems prevents binocular alignment. This leads to cortical suppression of the deviating eye in children to avoid diplopia, which may cause amblyopia if uncorrected.

\section*{Etiology}
\begin{itemize}
  \item \textbf{Infantile esotropia:} Large-angle inward deviation presenting before 6 months of age
  \item \textbf{Accommodative esotropia:} Convergence stimulated by uncorrected hyperopia
  \item \textbf{Intermittent exotropia:} Outward deviation, often worse with fatigue or distance viewing
  \item \textbf{Paralytic (cranial nerve palsy):} CN III (oculomotor), IV (trochlear), or VI (abducens) paresis
  \item \textbf{Restrictive:} Thyroid eye disease, orbital fracture, or prior retinal detachment surgery
  \item \textbf{Sensory:} Vision loss from cataract or corneal scar causing loss of fusion
\end{itemize}

\section*{Indications for Evaluation}
Seek care if:
\begin{itemize}
  \item Eye misalignment noted by parent or caregiver in a child of any age
  \item New-onset strabismus in an adult (suspect cranial nerve palsy or intracranial pathology)
  \item Diplopia (double vision) in any patient with known strabismus
  \item Head tilt or turn that the patient adopts to maintain binocularity
  \item Failed vision screening at school or pediatric examination
\end{itemize}

\section*{Related Symptoms}
\begin{itemize}
  \item Diplopia (more common in adult-onset strabismus)
  \item Amblyopia (in untreated childhood strabismus)
  \item Abnormal head posture (face turn, chin elevation/depression, head tilt)
  \item Impaired depth perception (reduced stereoacuity)
  \item Eye strain, headache, and visual fatigue (especially with intermittent deviations)
\end{itemize}
\textbf{Note:} In children, the angle of deviation and amblyopia status must be assessed before surgical planning; in adults with new-onset strabismus, neuroimaging is indicated to rule out compressive or ischemic cranial neuropathy.

\section*{Self-Care / Home Management}
\begin{itemize}
  \item Patching the stronger eye for prescribed hours to treat amblyopia in children
  \item Prism glasses may relieve diplopia while awaiting surgery
  \item Eye exercises (orthoptic therapy) for convergence insufficiency
  \item Avoiding precipitating factors in intermittent exotropia (fatigue, illness, bright sunlight)
  \item Sunglasses to reduce photophobia in intermittent exotropia (light sensitivity worsens the deviation)
\end{itemize}

\section*{Diagnostic Approach}
\begin{itemize}
  \item Cover-uncover test and alternate cover test for distance and near with and without correction
  \item Measurement of deviation using prism and cover test (Hirschberg and Krimsky methods in young children)
  \item Cycloplegic refraction to identify accommodative component
  \item Extraocular motility examination (versions, ductions) in nine cardinal positions of gaze
  \item Three-step test (Park's) for vertical deviations to identify the paretic muscle
  \item Hess chart or Lancaster red-green test for incomitant strabismus
  \item Neuroimaging for restrictive or paretic strabismus of unclear etiology
\end{itemize}

\section*{Treatment / Management Overview}
\begin{itemize}
  \item Optical correction (glasses) to manage accommodative esotropia
  \item Amblyopia therapy (patching or atropine penalization) before surgical intervention
  \item Prism glasses for diplopia relief in adults not candidates for surgery
  \item Botulinum toxin injection into the extraocular muscle as a diagnostic or therapeutic tool
  \item Strabismus surgery: Recession (weakening) and/or resection (strengthening) of appropriate muscles
  \item Adjustable suture technique for adults to fine-tune alignment postoperatively
  \item Orthoptic exercises for convergence insufficiency and intermittent exotropia
\end{itemize}

\clearpage
